%--- Ajustes del documento.
\pagestyle{plain}	% Páginas sólo con numeración inferior al pie

% -------------------------
%
% RESUMEN:
% 
%

% EDITAR: Resumen (máx. 1 pág.)
%\cleardoublepage % Se incluye para modificar el contador de página antes de añadir bookmark
\phantomsection  % Necesario con hyperref
\addcontentsline{toc}{chapter}{Resumen} % Añade al TOC.
\selectlanguage{spanish} % Selección de idioma del resumen.
\makeatletter
\begin{center} %
   {\textsc{TRABAJO FIN DE GRADO - ESCUELA SUPERIOR DE INFORMÁTICA 
   (UCLM)}\par} % Tipo de trabajo
   \vspace{1cm} %  
   {\textbf{\Large\@tituloCorto}\par}  % Título del trabajo
   \vspace{0.4cm} %
   {\@autor \\ \@cityTF,{} \@mesTF{} \@yearTF\par} 
   \vspace{0.9cm} %
   {\textbf{\large\textsf{Resumen}}\par} % Título de resumen
\end{center}   
\makeatother %
La codificación manual de diagnósticos clínicos en sistemas CIE-10 presenta importantes desafíos operativos: complejidad creciente (68,069 códigos vs 14,025 en CIE-9), tiempos promedios de 15 minutos por caso \cite{gogorcena2017cie10es}. 

Este trabajo propone un sistema automático de codificación mediante modelos BERT multi-etiqueta y con jerarquía, adaptados a narrativas clínicas en español, abordando tres problemáticas clave:

1) Procesamiento de lenguaje clínico no estructurado

2) Asignación múltiple de códigos jerárquicamente relacionados

3) Adaptación a requisitos regulatorios de la Orden SSI/1461/2022, normativa española emitida por el Ministerio de Sanidad.

La metodología combina técnicas avanzadas de PLN con conocimiento clínico estructurado:

1) Preprocesamiento específico para textos médicos (normalización de entidades clínicas y detección de negaciones)

2) Fine-tuning del modelo BETO (BERT para español) con capas de atención jerárquica

3) Mecanismo de pérdida adaptativa para manejar el desbalance extremo de clases. El sistema se entrena y evalúa usando los datos proporcionados en el Codiesp \cite{miranda2020}, un conjunto de informes codificados de forma manual por expertos, con datos de entrenamiento, test y validación.

Los resultados muestran un F1-score de //TODO \%// en codificación primaria, superando el rendimiento humano (0.89). El modelo reduce el tiempo de codificación //TODO, añadir datos// manteniendo compliance con niveles de especificidad requeridos.

El trabajo pretende demostrar la viabilidad de sistemas basados en transformers para codificación clínica automatizada, proponiendo un pipeline reproducible que integra: 

1) Un componente de clasificación automática basado en los datos de entrada

2) Un módulo de auditoría automática y gestión de incidencias

3) Una interfaz API REST para la integración con sistemas hospitalarios. 

Las limitaciones identificadas (sesgos en datos de entrenamiento, manejo de neologismos médicos) establecen directrices para futuras investigaciones.

\bigskip
\noindent\textbf{Palabras clave}: CIE-10, BERT, clasificación multilabel, codificación clínica, procesamiento de lenguaje natural.

%---
\cleardoublepage % Se incluye para modificar el contador de página antes de añadir 





% EDITAR: Abstract (máx. 1 pág.)
%---
\phantomsection  % Necesario con hyperref
\addcontentsline{toc}{chapter}{Abstract} % Añade al TOC.
\selectlanguage{english} % Selección de idioma del resumen.
\makeatletter
\begin{center} %
   {\textsc{BACHELOR DISSERTATION - ESCUELA SUPERIOR DE INFORMÁTICA 
   (UCLM)}\par}
   \vspace{1cm} %  
   {\textbf{\Large Guided template for TFG}\par}
   \vspace{0.4cm} %
   {\@autor \\ \@cityTF,{} \@monthTF{} \@yearTF\par} 
   \vspace{0.9cm} %
   {\textbf{\large\textsf{Abstract}}\par} 
\end{center}   
\makeatother %
\emph{English version for the abstract.}

\bigskip 

\noindent\textbf{Keywords}: \dots

\cleardoublepage % Se incluye para modificar el contador de página antes de añadir 

